A detecção de anomalias em mercados financeiros é dominada, na literatura, pelo
mercado norte-americano e por criptomoedas. Este trabalho aplica uma arquitetura
\emph{LSTM-Autoencoder} à detecção não supervisionada de anomalias em quatro ações
da B3 (PETR4, VALE3, AMER3 e ITUB4), cobrindo distintos setores econômicos. O modelo
é treinado apenas sobre o período considerado de ``normalidade'' (2010--2019) e sinaliza
como anômalas as janelas cujo erro de reconstrução excede um limiar; comparam-se
limiares estático (percentil fixo) e dinâmico (percentil em janela móvel causal).
A avaliação combina validação narrativa (sobreposição com eventos econômicos e políticos
brasileiros documentados) e quantitativa (injeção controlada de anomalias sintéticas,
com Precision, Recall e F1). Este relatório preliminar documenta, de forma incremental e
academicamente rastreável, as decisões metodológicas, suas fontes e as observações de
cada etapa do projeto. Até o momento concluíram-se as etapas de \emph{infraestrutura}
(M0) e de \emph{coleta e análise exploratória} (M1).
